\chapter{16}

\section{Donnerstag, 23. August}



«Noch n’ Tipp für dich. Bist ein schöner Kerl, gut für die Frauen,
schlecht für mich,» und Mascha lachte übertrieben ausgelassen, «kauf dir
eine neue Brille. Einfach Scheisse, dieses Ding in deinem hübschen
Gesicht.»

Damit liess Mascha Louis mit offenem Mund stehen. Er verschwand mit
hastigen Schritten.

\sectionbreak

Es hatte geklappt. Gegen alle Bedenken. War unkomplizierter abgelaufen
als gedacht, und Louis hielt nach zwei Tagen des Bangens ein neu
kreiertes Ich in den Händen. Die Strapazen hatten sich gelohnt. Hier
stand er. In dieser dunkeln Augustnacht, die sich ihm einbrennen würde.
Er starrte auf die Angaben zu seiner Person. Nahm die lästige Brille von
der Nase und rieb sich die Augen. Sie brannten wieder. Das klare Bild
der Umgebung wirkte wohltuend.

Unverschämt, wie gut das Dokument gemacht war, wie professionell es
aussah, dieses rote Büchlein. Ein völlig normaler Schweizer Pass. Ganz
selbstverständlich fügte sich sein Bild neben den Angaben ein. Er saugte
die Daten begierig in sich auf. Von nun an also hatte er die Wahl. Er
fühlte sich seltsam motiviert. Er konnte sich als Louis mit spanischem
Oder als Ciro mit Schweizer Pass ausgeben. Er war ab sofort mit sich zu
zweit. Die Stimme, die sich leise aus einem verborgenen Winkel meldete,
ignorierte er vehement. Nein, er war weder übergeschnappt noch dumm.
Einzig abgekämpft nach Tagen panischer Angst.

Louis senkte das Dokument und schaute stirnrunzelnd vor sich hin. Mascha
musste in der Szene tief verankert sein. Jess hatte recht. Er schien
Menschen einschätzen zu können. In jeder Hinsicht. Selbst Louis’ Grösse
stimmte annähernd. Und Mascha hatte tatsächlich seine Bitte erhört. Er
hatte einen spanischen Namen bekommen. Dass er ein gutes Jahr älter
ausgewiesen war, würde nicht auffallen. An einen Geburtstag im August
hingegen würde er sich gewöhnen müssen. Ja, Mascha war gut. Und er
fühlte sich ihm zugeneigt.

\sectionbreak

Louis steckte das Büchlein ein und schaute sich um. Es war Zeit, diesen
Ort zu verlassen.

Der Pass würde ihn vorübergehend schützen. Vor seiner eigenen
Geschichte.

\vspace{\baselineskip}

\textbf{Lorente Ciro José}

\textbf{Schweiz Suisse Svizzera Svizra Switzerland}

\textbf{14.08.1988  \hspace{4mm}  M  \hspace{8mm} 175cm}

Heimatort:

\textbf{Chêne-Bougeries GE}

Ausgestellt am:

\textbf{13.06.2014}

Gültig bis:

\textbf{12.06.2024} \hspace{1cm} Behörde Autorité: \textbf{Genf GE}

\vspace{\baselineskip}

Es war nach Mitternacht. Der wolkenlose Himmel hing dunkelgrau. Die
Müdigkeit sass Louis erneut schwer im Nacken. Und doch war es wohltuend,
einfach zu gehen. Er versuchte, sich auf das Auftreten seiner Fusssohlen
zu konzentrieren. Sein Kopf fühlte sich freier an. Nach einer Weile
erreichte er den Hauptbahnhof.

Für diese Nacht wollte er einen andern Schlafplatz finden.

Vom Europaplatz lenkte er sich der Sihl entlang Richtung
«Gessnerbrücke». Aus der Ferne fiel Louis ein Gebäude am Ende der Brücke
auf der andern Uferseite auf.

Er überquerte die Brücke und blieb vor dem kleinen Gebäude stehen. «Bar
Rio» las er oben auf einer länglichen Fensterfront. Louis studierte die
Umgebung. Auch um diese Zeit waren etliche Menschen unterwegs.

Er wartete, bis ein junges Paar ausser Sichtweite war. Über einen
lauschigen, kleinen Platz neben der Bar gelangte er zu einer
Eisentreppe. Sie führte hinunter zu dreistufig angelegten, hohen
Steinbänken. Die Dunkelheit strengte ihn an. Und teilweise tastend
erkundete er einen möglichen Weg. Dann fiel ihm die Taschenlampe ein.

Er leuchtete die Umgebung aus und stieg vorsichtig nach unten. Neben der
Treppe, auf oberster Stufe, befand er sich auf einem schmalen Platz.
Zweiseitig von dicht nebeneinander stehenden, niederen Bäumen umrahmt.
Der perfekte Ort. Er war erleichtert. Hier war er blickdicht geschützt.

Auf seiner Seite der Sihl und am gegenüberliegenden Ufer konnte Louis
keine Menschen ausmachen. Er hatte sein Nachtlager gefunden. Als finale
Tarnung käme nun seine dünne, grüne Decke zum Einsatz.
Auseinandersetzungen mit Ordnungshütern musste er ebenso vermeiden wie
Begegnungen mit Privatpersonen. Sein Energielevel stand auf null.

Beim Anbruch des Tages würde er verschwunden sein. Wie die beiden Nächte
zuvor.

Auch die hatte er schliesslich erfolgreich überstanden. Louis hatte
seinen Schlafsack in Seenähe, in gutem Abstand zu jenen ausgebreitet,
die ebenfalls draussen schliefen. Junge Leute hatten bis in die
Morgenstunden zusammen am Ufer des Zürichsees gesessen, Louis jeden
näheren Kontakt vermieden. Er hatte sich freundlich distanziert gegeben.
Wurde er angesprochen, erfand er irgendwelche Kurzgeschichten, leicht
variable, je nach Situation. Und wenn er sich in Gefahr sah, mimte er
den Franzosen, der kein Englisch sprach.

\sectionbreak

Louis richtete sich in der hintersten Ecke, direkt vor den Bäumen, ein.
Er rutschte hin und her, bis er eine einigermassen bequeme Stellung
fand. Der Untergrund war hart. Wieder kam die Sehnsucht nach einer
Dusche und einem frischen Bett auf. Er wünschte nichts mehr, als sich
für ein paar Stunden fallen zu lassen, abzuschalten und seine Sorgen zu
vergessen.

Damit fiel er in einen komaähnlichen Tiefschlaf.

\sectionbreak

Kaum erwacht, zog Louis am nächsten Morgen als Erstes seinen neuen Pass
aus der Seitentasche seiner Cargohose. Es war knapp sechs Uhr vorbei. Er
las die Angaben wieder und wieder. Von irgendwoher kam dennoch der
Gedanke, ob das Dokument einer scharfen Prüfung standhalten konnte.

Es half nichts. Eine Passkontrolle würde es zeigen. Es musste reichen,
zu wissen, dass es Totalfälschungen gab, Pässe, die komplett neu kreiert
wurden. Das konnte nicht möglich sein. In zwei Tagen nicht zu schaffen.
Es musste sich also um eine Verfälschung handeln. Dabei wurden nur
Passbilder oder Angaben wie Namen und Daten ausgetauscht.

Er musste seiner neuen Identität vertrauen. Unverzüglich. Er hatte das
Dokument. Basta.

Nach wie vor lag Louis in seinem angenehm warmen Schlafsack. Wenigstens
seine Neuanschaffungen bewährten sich. Diffus beflügelt begann er, sich
um «Ciro» eine passende Fassade aufzubauen. Ein sympathischer, starker
Typ wollte er sein. Das würde ihm im Kontakt mit neuen Menschen dienen.
Er könnte ein Musikerleben konstruieren. Was ihm leicht fiele. Er würde
überzeugen können. Und den Naturliebhaber geben. Zusätzlich vielleicht.
Seine Phantasien begannen ihm zu gefallen. Auch als Übersetzer könnte er
auftreten. Sprachaffin und ebenso geübt war er. Er würde seine
Kenntnisse anpreisen. Dabei huschte ein Lächeln über sein Gesicht. Wie
selbstverständlich er zwischen den verschiedenen Sprachen hin- und
herpendeln konnte.

Alles musste anders werden. Leicht würde es nicht. Aber womöglich
einfacher als bisher. Schluss war mit Erfinden von Namen und Geschichten
von jetzt auf gleich. In einer Geschwindigkeit, die ihm den Atem nahm.

Er würde aussteigen aus diesem Improvisationstheater. Würde seinen Text
kennen und üben können. Wegkommen von Hetze und Anstrengung. Sich neu
erschaffen. Bis er wusste, wie es weiterging.

\sectionbreak

Wie tags zuvor hatte Louis sich an Kiosken Titelseiten von
Tageszeitungen angesehen. Und war erleichtert. Keine Fahndung, keine
Zeugenaufrufe und vor allem, keine Bilder, weder von Giorgia noch von
ihm. Wiederum hatte er drei Tagesblätter gekauft. Nichts. Akribisch
genau hatte er jede Seite durchgesehen, auch kleinste Artikel
überflogen.

Später würde er seine Fahndungsangst geradezu als lächerlich abtun. Was
immer auch mit ihr geschehen war, so schnell wurde nicht gefahndet.

\sectionbreak

Louis suchte den Einstieg in die Bahnhofstrasse. Laut Fahrplan an einer
Tramhaltestelle endete sie auf dem Bürkliplatz. Nach ein paar Minuten
lag der Zürichsee vor ihm. Louis sah sich um, fand schliesslich ein
einladendes Restaurant am Seeufer und setzte sich an einen kleinen,
freien Tisch.

Beim Zuklappen der Speisekarte erschien das Bild. Überraschend und
ebenso willkommen. Er sah sich, wie er beglückt seine neue Gitarre aus
dem Koffer heben würde. Wie er sie kennenlernte. Und wie sich seine
Verlorenheit mit jedem Klang um eine Spur auflösen würde. Ein
verheissungsvoller Tag schien auf ihn zu warten. Er wollte darauf
hoffen.

Es dauerte, bis das Frühstück vor ihm stand. Und für einen flüchtigen
Moment fühlte er sich wie ein Urlauber.
